\documentclass[12pt,letterpaper]{article}

\usepackage[spanish,es-nodecimaldot]{babel}
\usepackage[utf8]{inputenc}
\usepackage[T1]{fontenc}
\usepackage{lmodern}
\usepackage[margin=1in]{geometry}
\usepackage{setspace}
\usepackage[hidelinks]{hyperref}
\usepackage{booktabs}
\usepackage{longtable}
\usepackage{tabularx}
\usepackage{array}
\usepackage{enumitem}
\usepackage{xcolor}
\usepackage{listings}
\usepackage{float}

\onehalfspacing
\setlength{\parindent}{0.5in}
\setlength{\parskip}{0pt}
\sloppy

\lstdefinestyle{pythonicociv}{
  language=Python,
  basicstyle=\ttfamily\small,
  keywordstyle=\color{blue!60!black},
  stringstyle=\color{green!40!black},
  commentstyle=\color{gray!70!black},
  showstringspaces=false,
  breaklines=true,
  frame=single,
  columns=fullflexible
}

\title{Documentación técnica de módulos de la aplicación ICOCIV}
\author{Esteban Ojeda}
\date{2026}

\begin{document}

\begin{titlepage}
    \centering
    \vspace*{1.5cm}
    {\Large\bfseries Documentación técnica de módulos de la aplicación ICOCIV\par}
    \vspace{0.5cm}
    {\large Reorganización de arquitectura, flujo de código y dependencias internas\par}
    \vspace{1.5cm}
    {\large Autor: Esteban Ojeda\par}
    {\large Universidad del Cauca\par}
    {\large Programa de Ingeniería Civil\par}
    \vfill
    {\large Popayán, Colombia\par}
    {\large 2026\par}
\end{titlepage}

\tableofcontents
\newpage

\section{Introducción}

Este documento describe la codificación y organización técnica de la aplicación ICOCIV después de la reorganización de módulos bajo una carpeta principal llamada \texttt{app\_icociv}. Su finalidad es permitir que una persona evaluadora, desarrolladora o revisora académica entienda el propósito de cada módulo sin tener que leer todo el código fuente de forma lineal.

La reorganización no modifica la metodología oficial del DANE ni altera la lógica estadística de la aplicación. Su objetivo es ordenar responsabilidades: la interfaz visual queda separada del backend, los módulos estadísticos quedan separados de la generación documental y las rutas de proyecto se centralizan para evitar dependencias frágiles.

\section{Propósito de la reorganización}

Antes de la reorganización, los módulos principales estaban distribuidos entre carpetas como \texttt{migracion/} e \texttt{interfaz/}. Aunque el backend funcionaba, la estructura no comunicaba con claridad qué archivos pertenecían a datos, estadística, modelos, reportes, persistencia o UI. La nueva estructura permite:

\begin{itemize}[leftmargin=*]
    \item agrupar todo el código activo en una sola carpeta principal;
    \item facilitar el mantenimiento académico y técnico;
    \item actualizar imports bajo un espacio de nombres claro;
    \item reducir ambigüedad entre validación de datos y validación predictiva;
    \item conservar datos, reportes, sesiones y documentación fuera del paquete de código.
\end{itemize}

\section{Estructura general de carpetas}

\begin{lstlisting}[style=pythonicociv]
app_icociv/
|-- config/
|-- datos/
|-- estadistica/
|-- modelos/
|-- validacion/
|-- proyeccion/
|-- reportes/
|-- exportables/
|-- persistencia/
|-- servicios/
|-- interfaz/
`-- utilidades/
\end{lstlisting}

La carpeta \texttt{app\_icociv} contiene únicamente módulos de aplicación. Los anexos Excel, documentos académicos, sesiones JSON y reportes generados permanecen en carpetas externas para no mezclar código fuente con insumos o resultados.

\section{Descripción de módulos de configuración}

\begin{longtable}{p{0.18\textwidth}p{0.23\textwidth}p{0.27\textwidth}p{0.16\textwidth}p{0.16\textwidth}}
\caption{Módulos de configuración}
\label{tab:modulos_configuracion}\\
\toprule
Módulo & Ubicación & Propósito & Entradas & Salidas \\
\midrule
\endfirsthead
\toprule
Módulo & Ubicación & Propósito & Entradas & Salidas \\
\midrule
\endhead
\texttt{rutas.py} & \texttt{app\_icociv/config/rutas.py} & Centraliza la raíz del proyecto y carpetas de sesiones, reportes y documentación. & Ubicación física del archivo. & Objetos \texttt{Path}. \\
\texttt{settings.py} & \texttt{app\_icociv/config/settings.py} & Define nombre y versión de la aplicación. & Ninguna. & Constantes de identificación. \\
\texttt{constantes.py} & \texttt{app\_icociv/config/constantes.py} & Define extensiones soportadas y horizontes UI predeterminados. & Ninguna. & Tuplas de constantes. \\
\bottomrule
\end{longtable}

\section{Descripción de módulos de datos}

\begin{longtable}{p{0.18\textwidth}p{0.23\textwidth}p{0.27\textwidth}p{0.16\textwidth}p{0.16\textwidth}}
\caption{Módulos de datos}
\label{tab:modulos_datos}\\
\toprule
Módulo & Ubicación & Propósito & Entradas & Salidas \\
\midrule
\endfirsthead
\toprule
Módulo & Ubicación & Propósito & Entradas & Salidas \\
\midrule
\endhead
\texttt{cargador\_datos.py} & \texttt{app\_icociv/datos/cargador\_datos.py} & Lee anexos oficiales Excel o XLSB, detecta encabezados reales y extrae tablas A16.x. & Bytes del archivo y nombre. & Diccionario de tablas y periodos detectados. \\
\bottomrule
\end{longtable}

Este módulo es la puerta de entrada de datos. No debe contener reglas de proyección ni lógica visual. Si se agregan nuevos anexos o formatos, la ampliación debe ocurrir aquí o en un módulo vecino dentro de \texttt{datos/}.

\section{Descripción de módulos estadísticos}

\begin{longtable}{p{0.17\textwidth}p{0.22\textwidth}p{0.27\textwidth}p{0.17\textwidth}p{0.17\textwidth}}
\caption{Descripción técnica de módulos estadísticos}
\label{tab:modulos_estadisticos}\\
\toprule
Módulo & Ubicación & Propósito & Entradas & Salidas \\
\midrule
\endfirsthead
\toprule
Módulo & Ubicación & Propósito & Entradas & Salidas \\
\midrule
\endhead
\texttt{analisis\_series.py} & \texttt{app\_icociv/estadistica/} & Normaliza periodos, valida serie mensual, calcula variables derivadas y factibilidad descriptiva. & Serie con periodo e índice. & Validaciones, derivadas y mensajes. \\
\texttt{criterios.py} & \texttt{app\_icociv/estadistica/} & Centraliza umbrales, criterios internos y matriz de auditoría metodológica. & Constantes y criterios declarados. & Matriz exportable y valores de referencia. \\
\texttt{diagnostico\_residuos.py} & \texttt{app\_icociv/estadistica/} & Calcula Durbin-Watson, ACF, PACF, Ljung-Box opcional y advertencias residuales. & Residuos del modelo. & Diagnóstico residual estructurado. \\
\texttt{metricas.py} & \texttt{app\_icociv/estadistica/} & Calcula MAE, RMSE, MAPE, sMAPE, MASE, AIC y AICc. & Observados, predichos y entrenamiento. & Métricas numéricas. \\
\texttt{modelos\_interpretables.py} & \texttt{app\_icociv/estadistica/} & Ajusta modelos base, Holt lineal, Holt amortiguado, regresiones y modelos sobre variación cuando aplican. & Tiempo e índice. & Modelos candidatos, parámetros y predicciones. \\
\texttt{validacion\_series.py} & \texttt{app\_icociv/estadistica/} & Realiza validaciones complementarias de serie temporal y outliers extremos. & Serie temporal. & Resumen de validación y alertas. \\
\bottomrule
\end{longtable}

Los módulos estadísticos no deben importar clases de PySide6 ni escribir archivos de reporte. Su salida debe ser información estructurada para que la capa de proyección o reportes decida cómo presentarla.

\section{Descripción de modelos predictivos}

\begin{longtable}{p{0.18\textwidth}p{0.23\textwidth}p{0.27\textwidth}p{0.16\textwidth}p{0.16\textwidth}}
\caption{Módulos de modelos predictivos}
\label{tab:modulos_modelos}\\
\toprule
Módulo & Ubicación & Propósito & Entradas & Salidas \\
\midrule
\endfirsthead
\toprule
Módulo & Ubicación & Propósito & Entradas & Salidas \\
\midrule
\endhead
\texttt{regresion.py} & \texttt{app\_icociv/modelos/regresion.py} & Contiene regresiones tradicionales y cálculo de intervalos bootstrap heredados. & Vectores de tiempo e índice. & Mejor modelo y parámetros. \\
\bottomrule
\end{longtable}

El módulo \texttt{modelos/regresion.py} se conserva porque el backend lo usa para compatibilidad con rutinas previas. La comparación estadística moderna se coordina desde \texttt{estadistica/modelos\_interpretables.py} y \texttt{proyeccion/servicio\_proyeccion.py}.

\section{Descripción de módulos de validación}

\begin{longtable}{p{0.18\textwidth}p{0.23\textwidth}p{0.27\textwidth}p{0.16\textwidth}p{0.16\textwidth}}
\caption{Módulos de validación predictiva}
\label{tab:modulos_validacion}\\
\toprule
Módulo & Ubicación & Propósito & Entradas & Salidas \\
\midrule
\endfirsthead
\toprule
Módulo & Ubicación & Propósito & Entradas & Salidas \\
\midrule
\endhead
\texttt{backtesting.py} & \texttt{app\_icociv/validacion/backtesting.py} & Ejecuta walk-forward validation, compara modelos y resume estabilidad predictiva. & Serie mensual, horizonte y modelo. & Predicciones históricas, errores y métricas. \\
\bottomrule
\end{longtable}

Este módulo representa la validación fuera de muestra. No debe leer Excel ni generar reportes; recibe una serie ya preparada y devuelve evidencia predictiva.

\section{Descripción de módulos de proyección}

\begin{longtable}{p{0.18\textwidth}p{0.23\textwidth}p{0.27\textwidth}p{0.16\textwidth}p{0.16\textwidth}}
\caption{Módulos de proyección}
\label{tab:modulos_proyeccion}\\
\toprule
Módulo & Ubicación & Propósito & Entradas & Salidas \\
\midrule
\endfirsthead
\toprule
Módulo & Ubicación & Propósito & Entradas & Salidas \\
\midrule
\endhead
\texttt{servicio\_proyeccion.py} & \texttt{app\_icociv/proyeccion/} & Resuelve selección jerárquica, construye serie, evalúa factibilidad, intervalos, horizontes y proyección final. & Tablas, selección, serie y horizonte solicitado. & Resultado integral de análisis y proyección. \\
\bottomrule
\end{longtable}

Este módulo es el núcleo analítico del backend. Aunque coordina varias etapas, no debe incorporar UI ni abrir archivos desde diálogos. La interfaz o el orquestador entregan los insumos.

\section{Descripción de módulos de reportes y exportables}

\begin{longtable}{p{0.18\textwidth}p{0.23\textwidth}p{0.27\textwidth}p{0.16\textwidth}p{0.16\textwidth}}
\caption{Módulos de reportes y exportables}
\label{tab:modulos_reportes_exportables}\\
\toprule
Módulo & Ubicación & Propósito & Entradas & Salidas \\
\midrule
\endfirsthead
\toprule
Módulo & Ubicación & Propósito & Entradas & Salidas \\
\midrule
\endhead
\texttt{generador\_reportes.py} & \texttt{app\_icociv/reportes/} & Genera PDF, DOCX, HTML, gráficas, tablas y textos técnicos. & Serie, resultado de proyección, selección y trazabilidad. & Archivos PDF, DOCX, HTML y figuras. \\
\texttt{csv\_reproducible.py} & \texttt{app\_icociv/exportables/} & Expone la exportación CSV reproducible como módulo funcional. & Serie, proyección y ruta jerárquica. & CSV con observado, proyectado, intervalo del 95\,\% y metadatos. \\
\bottomrule
\end{longtable}

Los reportes presentan resultados; no deben recalcular el modelo. Si se detecta una diferencia entre informe y backend, debe corregirse el backend o la extracción de datos del resultado, no duplicar fórmulas en el reporte.

\section{Descripción de persistencia, utilidades e interfaz}

\begin{longtable}{p{0.18\textwidth}p{0.23\textwidth}p{0.27\textwidth}p{0.16\textwidth}p{0.16\textwidth}}
\caption{Módulos de soporte, persistencia e interfaz}
\label{tab:modulos_soporte_interfaz}\\
\toprule
Módulo & Ubicación & Propósito & Entradas & Salidas \\
\midrule
\endfirsthead
\toprule
Módulo & Ubicación & Propósito & Entradas & Salidas \\
\midrule
\endhead
\texttt{gestor\_sesiones.py} & \texttt{app\_icociv/persistencia/} & Guarda, carga y lista sesiones JSON locales. & Datos de sesión y ruta. & Archivos JSON y diccionarios. \\
\texttt{utilidades.py} & \texttt{app\_icociv/utilidades/} & Normaliza periodos, filtra DataFrames y calcula auxiliares estadísticos. & Textos, DataFrames o arrays. & Valores normalizados o filtrados. \\
\texttt{nomenclatura\_icociv.py} & \texttt{app\_icociv/utilidades/} & Centraliza nombres amigables de tablas, niveles y ruta jerárquica. & Claves técnicas ICOCIV. & Textos visibles y rutas normalizadas. \\
\texttt{ventana\_principal.py} & \texttt{app\_icociv/interfaz/} & Construye la ventana PySide6, gráficos, selección de archivo, sesiones y botones de reporte. & Interacción del usuario. & UI y llamadas a controladores. \\
\texttt{controlador\_principal.py} & \texttt{app\_icociv/interfaz/controladores/} & Coordina selección jerárquica y llamadas al backend desde UI. & Tablas, selección y eventos. & Estado de UI y resultados. \\
\texttt{trabajadores.py} & \texttt{app\_icociv/interfaz/controladores/} & Ejecuta funciones largas en hilos Qt. & Función y argumentos. & Señales de resultado o error. \\
\texttt{modelo\_tabla.py} & \texttt{app\_icociv/interfaz/widgets/} & Adapta DataFrames a tablas Qt. & DataFrame. & Modelo Qt para tabla. \\
\bottomrule
\end{longtable}

\section{Descripción de servicios y orquestación}

\begin{longtable}{p{0.18\textwidth}p{0.23\textwidth}p{0.27\textwidth}p{0.16\textwidth}p{0.16\textwidth}}
\caption{Módulos de servicios y orquestación}
\label{tab:modulos_servicios}\\
\toprule
Módulo & Ubicación & Propósito & Entradas & Salidas \\
\midrule
\endfirsthead
\toprule
Módulo & Ubicación & Propósito & Entradas & Salidas \\
\midrule
\endhead
\texttt{orquestador.py} & \texttt{app\_icociv/servicios/} & Coordina carga de archivo, análisis completo y proyección de serie sin duplicar lógica interna. & Ruta de archivo, selección y horizonte. & Resultado integral del backend. \\
\bottomrule
\end{longtable}

El orquestador es una fachada técnica. Su función es documentar el flujo y ofrecer un punto de entrada para scripts o pruebas, sin mezclar funciones estadísticas de bajo nivel.

\section{Flujo general de ejecución}

\begin{enumerate}[leftmargin=*]
    \item El usuario abre la app con \texttt{python aplicacion.py}.
    \item La interfaz permite seleccionar un archivo Excel o XLSB.
    \item \texttt{datos/cargador\_datos.py} lee tablas y periodos.
    \item El controlador actualiza los desplegables jerárquicos.
    \item \texttt{proyeccion/servicio\_proyeccion.py} resuelve la fila seleccionada.
    \item Se construye la serie mensual con periodo e índice.
    \item \texttt{estadistica/analisis\_series.py} valida, normaliza y calcula variables derivadas.
    \item \texttt{validacion/backtesting.py} ejecuta walk-forward.
    \item \texttt{estadistica/metricas.py} calcula errores y métricas.
    \item \texttt{proyeccion/servicio\_proyeccion.py} clasifica horizontes e intervalos.
    \item \texttt{reportes/generador\_reportes.py} exporta PDF, DOCX, HTML o CSV.
\end{enumerate}

\section{Dependencias internas entre módulos}

\begin{longtable}{p{0.28\textwidth}p{0.32\textwidth}p{0.32\textwidth}}
\caption{Relación entre capas internas}
\label{tab:dependencias_internas}\\
\toprule
Capa & Depende de & Usada por \\
\midrule
\endfirsthead
\toprule
Capa & Depende de & Usada por \\
\midrule
\endhead
\texttt{datos} & \texttt{utilidades} & \texttt{interfaz}, \texttt{servicios}, \texttt{proyeccion} \\
\texttt{estadistica} & \texttt{utilidades} & \texttt{validacion}, \texttt{proyeccion}, \texttt{reportes} \\
\texttt{modelos} & \texttt{utilidades}, librerías científicas & \texttt{proyeccion} \\
\texttt{validacion} & \texttt{estadistica}, \texttt{modelos} & \texttt{proyeccion} \\
\texttt{proyeccion} & \texttt{datos}, \texttt{estadistica}, \texttt{validacion}, \texttt{modelos} & \texttt{interfaz}, \texttt{servicios}, \texttt{reportes} \\
\texttt{reportes} & \texttt{estadistica}, \texttt{persistencia}, \texttt{utilidades} & \texttt{interfaz}, \texttt{servicios} \\
\texttt{interfaz} & \texttt{config}, \texttt{datos}, \texttt{proyeccion}, \texttt{reportes}, \texttt{persistencia} & Usuario final \\
\bottomrule
\end{longtable}

\section{Guía para agregar nuevos modelos}

Para agregar un modelo:

\begin{enumerate}[leftmargin=*]
    \item Implementar el ajuste en \texttt{estadistica/modelos\_interpretables.py} o en \texttt{modelos/} si es una utilidad independiente.
    \item Registrar nombre, tipo y razón metodológica en el catálogo de \texttt{proyeccion/servicio\_proyeccion.py}.
    \item Asegurar que el modelo pueda ser reentrenado en cada origen walk-forward.
    \item Definir parámetros reproducibles para reportes.
    \item Agregar pruebas de importación, backtesting y reporte.
\end{enumerate}

\section{Guía para agregar nuevos criterios}

Los criterios deben registrarse en \texttt{estadistica/criterios.py}. Cada criterio nuevo debe indicar:

\begin{itemize}[leftmargin=*]
    \item nombre técnico;
    \item tipo de sustento;
    \item valor o regla;
    \item impacto en la decisión;
    \item módulos donde se usa;
    \item si puede recalibrarse.
\end{itemize}

\section{Guía para modificar reportes}

Los reportes deben leer resultados ya calculados. No deben volver a ajustar modelos ni cambiar la factibilidad. Si una sección del informe necesita datos nuevos, primero se deben agregar al resultado del backend y luego presentarlos en \texttt{generador\_reportes.py}.

\section{Buenas prácticas de mantenimiento}

\begin{itemize}[leftmargin=*]
    \item Mantener imports bajo \texttt{app\_icociv.*}.
    \item Evitar rutas relativas basadas en la ubicación física de un widget.
    \item Usar \texttt{config/rutas.py} para carpetas del proyecto.
    \item Mantener docstrings en español para funciones principales.
    \item No mezclar lógica estadística con interfaz.
    \item No duplicar fórmulas en reportes.
    \item Verificar \texttt{python tests/test\_imports\_modulos.py} después de cambios de arquitectura.
\end{itemize}

\section{Conclusiones}

La reorganización bajo \texttt{app\_icociv/} permite que la aplicación se entienda como un sistema modular: datos, estadística, modelos, validación, proyección, reportes, exportables e interfaz tienen responsabilidades separadas. Esta separación mejora la mantenibilidad, facilita la revisión académica y reduce el riesgo de romper el backend estadístico al modificar la UI o los informes.

\end{document}
